% answers_short.tex —— 2006 年·普通高等学校招生考试·江西卷理科 · 简短答案速查
% 由 tools/make_short_answers.py 自 answers.tex 的 \gksolhead/\gkanswer 提取，
% 并与 2006/2006_数学真题_江西理_答案要点.md 逐题比对后生成，勿手改（改源 MD 后重跑）。

\begin{gknotice}
\textbf{使用说明}：本节只给最终答案便于快速核对，完整推导见解析版。答案为编写者独立推导，\textbf{非官方标准答案}。
\end{gknotice}

\gkpart{选择题}
\begin{center}\begin{tabular}{ccccccccccccc}
\toprule
题号 & 1 & 2 & 3 & 4 & 5 & 6 & 7 & 8 & 9 & 10 & 11 & 12 \\
\midrule
答案 & C & D & D & B & C & C & D & B & D & A & C & A \\
\bottomrule
\end{tabular}\end{center}

\gkpart{填空题}
\begin{itemize}[leftmargin=2.6em]
  \item[\textbf{13.}] $1/2$。
  \item[\textbf{14.}] $2$。
  \item[\textbf{15.}] $5\sqrt2$。
  \item[\textbf{16.}] ②④。
\end{itemize}

\gkpart{解答题}
\begin{itemize}[leftmargin=2.6em]
  \item[\textbf{17.}] (1) $a=-1/2,\ b=-2$；递增区间 $(-\infty,-2/3)$、$(1,+\infty)$，递减区间 $(-2/3,1)$；(2) $c<-1$ 或 $c>2$
  \item[\textbf{18.}] (1) $\xi$ 取 $0,10,20,50,60$ 的概率依次为 $729/1000,243/1000,18/1000,9/1000,1/1000$；(2) $E(\xi)=33/10$ 元
  \item[\textbf{19.}] (1) $S_1=\frac{\sin\alpha}{12\sin(150^\circ-\alpha)}$，$S_2=\frac{\sin\alpha}{12\sin(\alpha-30^\circ)}$；(2) 最大值 $240$，最小值 $216$
  \item[\textbf{20.}] (1) $AD\perp BC$；(2) 二面角 $B-AC-D$ 为 $\arccos\sqrt{2/3}$；(3) 存在两点；若 $E=A+t(C-A)$，则 $t=1\pm\sqrt2/2$
  \item[\textbf{21.}] (1) 点 $P$ 的轨迹方程为 $\frac{x^2}{a^2}+\frac{y^2}{b^2}=\frac{cx}{a^2}$；(2) $\theta=\pi/2$；此时 $m$ 为 $x=1$，三角形面积最大值为 $1/\sqrt2$
  \item[\textbf{22.}] (1) $a_n=\frac{n}{1-3^{-n}}=\frac{n3^n}{3^n-1}$；(2) 见解析版
\end{itemize}
